% ---------------------------------------------------------
% TABELA EMENTA 31: Oficina de Integração I — Semestre 4 e 5
% ---------------------------------------------------------
\begin{xltabular}{\linewidth}{|X|p{2.5cm}|p{2.5cm}|}
\hline
\multirow{3}{=}{\textcolor{ifscgreen}{\textbf{Unidade Curricular:}}\\[3pt]\textbf{\large Oficina de Integração I — Semestre 4 e 5}} & \multicolumn{2}{p{\dimexpr5cm+2\tabcolsep+\arrayrulewidth\relax}|}{\textcolor{ifscgreen}{\textbf{Semestre:}} \textbf{4º e 5º (transição Ano 2/Ano 3)}} \\ \cline{2-3}
 & \textcolor{ifscgreen}{\textbf{CH EaD*:}} & \textcolor{ifscgreen}{\textbf{CH Total*:}} \\ \cline{2-3}
 & \textbf{00 h} & \textbf{80 h} \\ \hline
\endhead
\multicolumn{3}{|p{\dimexpr\linewidth-2\tabcolsep-2\arrayrulewidth\relax}|}{\textcolor{ifscgreen}{\textbf{Objetivos:}}} \\ \hline
\multicolumn{3}{|p{\dimexpr\linewidth-2\tabcolsep-2\arrayrulewidth\relax}|}{
\begin{itemize}[leftmargin=*,noitemsep,topsep=2pt]
 \item Integrar conhecimentos da formação geral e da formação técnica na compreensão de fenômenos relacionados às organizações, ao trabalho, à ciência e ao território.
 \item Estabelecer as múltiplas relações entre as diferentes áreas do conhecimento por meio de atividades teóricas e práticas associadas às temáticas trabalhadas.
 \item Desenvolver competências investigativas por meio da observação, problematização e análise de situações reais.
 \item Mobilizar conhecimentos científicos, técnicos e culturais na construção de interpretações sobre as temáticas estudadas.
 \item Participar de ciclos temáticos interdisciplinares que promovam a articulação entre ensino, pesquisa e extensão.
 \item Produzir sínteses individuais e coletivas, utilizando diferentes linguagens e formas de comunicação científica e técnica.
 \item Desenvolver materiais concretos, produções escritas e audiovisuais relacionados aos temas e temáticas da unidade curricular.
\end{itemize}
} \\ \hline
\multicolumn{3}{|p{\dimexpr\linewidth-2\tabcolsep-2\arrayrulewidth\relax}|}{\textcolor{ifscgreen}{\textbf{Conteúdos:}}} \\ \hline
\multicolumn{3}{|p{\dimexpr\linewidth-2\tabcolsep-2\arrayrulewidth\relax}|}{
\begin{itemize}[leftmargin=*,noitemsep,topsep=2pt]
 \item Integração curricular e interdisciplinaridade.
 \item Ciência como princípio educativo.
 \item Métodos de investigação e produção do conhecimento.
 \item Observação, diagnóstico e análise da realidade organizacional e territorial.
 \item Comunicação científica e técnica.
 \item Produção de sínteses integradoras.
 \item Demais conteúdos envolvidos nas temáticas da oficina.
\end{itemize}
} \\ \hline
\multicolumn{3}{|p{\dimexpr\linewidth-2\tabcolsep-2\arrayrulewidth\relax}|}{\textcolor{ifscgreen}{\textbf{Estratégias de Ensino e Aprendizagem:}}} \\ \hline
\multicolumn{3}{|p{\dimexpr\linewidth-2\tabcolsep-2\arrayrulewidth\relax}|}{- **Metodologia:** Componente de natureza interdisciplinar organizado em ciclos temáticos, definidos coletivamente pela equipe docente, articulando componentes da formação técnica e da formação geral. Cada ciclo é desenvolvido em etapas (segundo Silva, 2016): Problematização, Instrumentalização, Experimentação (oficinas, estudos de caso, pesquisas, visitas técnicas, observações de campo), Orientação, Sistematização, Consolidação (relatórios, painéis, portfólios, apresentações) e Socialização dos conhecimentos produzidos. Estratégias: aulas expositivas-dialogadas, análise documental, metodologias ativas (aprendizagem baseada em problemas, estudos de caso, investigação científica), oficinas, seminários, visitas técnicas, pesquisas de campo, rodas de conversa e laboratórios de aprendizagem, com atuação integrada de docentes da formação profissional e técnica. Ambientes: salas de aula, biblioteca, laboratório de informática, Laboratório de Inovação, Empreendedorismo e Desenvolvimento Sustentável (LINEDS), além de organizações e espaços comunitários.
- **Avaliação:** Diagnóstica, processual e formativa, considerando a participação dos estudantes, o desenvolvimento das atividades investigativas, a capacidade de integração dos conhecimentos, a produção de sínteses individuais e coletivas, a comunicação dos resultados e a reflexão crítica sobre os processos desenvolvidos. Instrumentos: relatórios, diários de bordo, portfólios, seminários, painéis, infográficos, mapas conceituais, estudos de caso, apresentações orais, autoavaliação e avaliação por pares.} \\ \hline
\multicolumn{3}{|p{\dimexpr\linewidth-2\tabcolsep-2\arrayrulewidth\relax}|}{\textcolor{ifscgreen}{\textbf{Bibliografia Básica:}}} \\ \hline
\multicolumn{3}{|p{\dimexpr\linewidth-2\tabcolsep-2\arrayrulewidth\relax}|}{1. FRIGOTTO, Gaudêncio; CIAVATTA, Maria; RAMOS, Marise Nogueira (org.). \textbf{Ensino médio integrado: concepção e contradições}. 3. ed. São Paulo: Cortez, 2012.
2. GIL, Antonio Carlos. \textbf{Como elaborar projetos de pesquisa}. 6. ed. São Paulo: Atlas, 2017.
3. MORIN, Edgar. \textbf{Os sete saberes necessários à educação do futuro}. 2. ed. rev. São Paulo: Cortez, 2011.
4. SILVA, Adriano Larentes da (org.) et al. \textbf{O currículo integrado no cotidiano da sala de aula}. Florianópolis: Publicação do IFSC, 2016.} \\ \hline
\multicolumn{3}{|p{\dimexpr\linewidth-2\tabcolsep-2\arrayrulewidth\relax}|}{\textcolor{ifscgreen}{\textbf{Bibliografia Complementar:}}} \\ \hline
\multicolumn{3}{|p{\dimexpr\linewidth-2\tabcolsep-2\arrayrulewidth\relax}|}{1. KEELLING, Ralph; BRANCO, Renato Henrique Ferreira. \textbf{Gestão de projetos: uma abordagem global}. 4. ed. São Paulo: Saraiva, 2019.
2. MARCONI, Marina de Andrade; LAKATOS, Eva Maria. \textbf{Metodologia do trabalho científico: projetos de pesquisa, pesquisa bibliográfica, teses de doutorado, dissertações de mestrado, trabalhos de conclusão de curso}. 8. ed. São Paulo: Atlas, 2018.} \\ \hline
\end{xltabular}

\vspace{0.4cm}


% ---------------------------------------------------------
% TABELA EMENTA 32: Oficina de Integração I — Ano 2
% ---------------------------------------------------------
\begin{xltabular}{\linewidth}{|X|p{2.5cm}|p{2.5cm}|}
\hline
\multirow{3}{=}{\textcolor{ifscgreen}{\textbf{Unidade Curricular:}}\\[3pt]\textbf{\large Oficina de Integração I — Ano 2}} & \multicolumn{2}{p{\dimexpr5cm+2\tabcolsep+\arrayrulewidth\relax}|}{\textcolor{ifscgreen}{\textbf{Semestre:}} \textbf{3º e 4º}} \\ \cline{2-3}
 & \textcolor{ifscgreen}{\textbf{CH EaD*:}} & \textcolor{ifscgreen}{\textbf{CH Total*:}} \\ \cline{2-3}
 & \textbf{00 h} & \textbf{80 h} \\ \hline
\endhead
\multicolumn{3}{|p{\dimexpr\linewidth-2\tabcolsep-2\arrayrulewidth\relax}|}{\textcolor{ifscgreen}{\textbf{Objetivos:}}} \\ \hline
\multicolumn{3}{|p{\dimexpr\linewidth-2\tabcolsep-2\arrayrulewidth\relax}|}{
\begin{itemize}[leftmargin=*,noitemsep,topsep=2pt]
 \item \textit{[Objetivos de aprendizagem pendentes de preenchimento pelo docente responsável]}
\end{itemize}
} \\ \hline
\multicolumn{3}{|p{\dimexpr\linewidth-2\tabcolsep-2\arrayrulewidth\relax}|}{\textcolor{ifscgreen}{\textbf{Conteúdos:}}} \\ \hline
\multicolumn{3}{|p{\dimexpr\linewidth-2\tabcolsep-2\arrayrulewidth\relax}|}{
\begin{itemize}[leftmargin=*,noitemsep,topsep=2pt]
 \item \textit{[Conteúdo programático pendente de preenchimento pelo docente responsável]}
\end{itemize}
} \\ \hline
\multicolumn{3}{|p{\dimexpr\linewidth-2\tabcolsep-2\arrayrulewidth\relax}|}{\textcolor{ifscgreen}{\textbf{Estratégias de Ensino e Aprendizagem:}}} \\ \hline
\multicolumn{3}{|p{\dimexpr\linewidth-2\tabcolsep-2\arrayrulewidth\relax}|}{\textit{[Metodologia de ensino e avaliação pendente de preenchimento pelo docente responsável]}} \\ \hline
\multicolumn{3}{|p{\dimexpr\linewidth-2\tabcolsep-2\arrayrulewidth\relax}|}{\textcolor{ifscgreen}{\textbf{Bibliografia Básica:}}} \\ \hline
\multicolumn{3}{|p{\dimexpr\linewidth-2\tabcolsep-2\arrayrulewidth\relax}|}{1. \textit{[1º título da Bibliografia Básica pendente de indicação docente e validação no Parecer da Biblioteca do Câmpus]}
2. \textit{[2º título da Bibliografia Básica pendente de indicação docente e validação no Parecer da Biblioteca do Câmpus]}} \\ \hline
\multicolumn{3}{|p{\dimexpr\linewidth-2\tabcolsep-2\arrayrulewidth\relax}|}{\textcolor{ifscgreen}{\textbf{Bibliografia Complementar:}}} \\ \hline
\multicolumn{3}{|p{\dimexpr\linewidth-2\tabcolsep-2\arrayrulewidth\relax}|}{1. \textit{[1º título da Bibliografia Complementar pendente de indicação docente e validação no Parecer da Biblioteca do Câmpus]}
2. \textit{[2º título da Bibliografia Complementar pendente de indicação docente e validação no Parecer da Biblioteca do Câmpus]}
3. \textit{[3º título da Bibliografia Complementar pendente de indicação docente e validação no Parecer da Biblioteca do Câmpus]}} \\ \hline
\end{xltabular}

\vspace{0.4cm}


% ---------------------------------------------------------
% TABELA EMENTA 33: Oficina de Integração II
% ---------------------------------------------------------
\begin{xltabular}{\linewidth}{|X|p{2.5cm}|p{2.5cm}|}
\hline
\multirow{3}{=}{\textcolor{ifscgreen}{\textbf{Unidade Curricular:}}\\[3pt]\textbf{\large Oficina de Integração II}} & \multicolumn{2}{p{\dimexpr5cm+2\tabcolsep+\arrayrulewidth\relax}|}{\textcolor{ifscgreen}{\textbf{Semestre:}} \textbf{5º e 6º (Ano 3)}} \\ \cline{2-3}
 & \textcolor{ifscgreen}{\textbf{CH EaD*:}} & \textcolor{ifscgreen}{\textbf{CH Total*:}} \\ \cline{2-3}
 & \textbf{00 h} & \textbf{80 h} \\ \hline
\endhead
\multicolumn{3}{|p{\dimexpr\linewidth-2\tabcolsep-2\arrayrulewidth\relax}|}{\textcolor{ifscgreen}{\textbf{Objetivos:}}} \\ \hline
\multicolumn{3}{|p{\dimexpr\linewidth-2\tabcolsep-2\arrayrulewidth\relax}|}{
\begin{itemize}[leftmargin=*,noitemsep,topsep=2pt]
 \item Integrar conhecimentos da formação profissional e formação geral para compreensão de problemas e desenvolvimento de propostas de solução.
 \item Estabelecer as múltiplas relações entre as diferentes áreas do conhecimento por meio de atividades teóricas e práticas associadas às temáticas trabalhadas.
 \item Planejar e desenvolver, de forma colaborativa, uma proposta de intervenção ou solução para um problema identificado.
 \item Comunicar o processo de desenvolvimento da proposição e os resultados obtidos, utilizando diferentes formas de apresentação.
 \item Desenvolver materiais concretos, produções escritas e audiovisuais relacionados aos temas e temáticas da unidade curricular.
\end{itemize}
} \\ \hline
\multicolumn{3}{|p{\dimexpr\linewidth-2\tabcolsep-2\arrayrulewidth\relax}|}{\textcolor{ifscgreen}{\textbf{Conteúdos:}}} \\ \hline
\multicolumn{3}{|p{\dimexpr\linewidth-2\tabcolsep-2\arrayrulewidth\relax}|}{
\begin{itemize}[leftmargin=*,noitemsep,topsep=2pt]
 \item Integração curricular aplicada.
 \item Planejamento e implementação de intervenções.
 \item Pesquisa aplicada para subsidiar a tomada de decisão.
 \item Desenvolvimento de soluções organizacionais e sociais.
 \item Gestão da execução.
 \item Monitoramento e avaliação de resultados.
 \item Comunicação científica e técnica.
 \item Socialização das aprendizagens.
\end{itemize}
} \\ \hline
\multicolumn{3}{|p{\dimexpr\linewidth-2\tabcolsep-2\arrayrulewidth\relax}|}{\textcolor{ifscgreen}{\textbf{Estratégias de Ensino e Aprendizagem:}}} \\ \hline
\multicolumn{3}{|p{\dimexpr\linewidth-2\tabcolsep-2\arrayrulewidth\relax}|}{- **Metodologia:** A unidade curricular é organizada em ciclos temáticos de intervenção definidos coletivamente pela equipe docente (formação geral e profissional, com possível participação de parceiros externos). Cada ciclo segue um percurso metodológico de sete momentos (inspirado em Silva, 2017): (1) Problematização; (2) Planejamento e Instrumentalização; (3) Implementação; (4) Orientação; (5) Monitoramento e Avaliação; (6) Consolidação; (7) Socialização. São utilizadas metodologias ativas como Design Thinking, Aprendizagem Baseada em Problemas, Aprendizagem Baseada em Projetos, estudos de caso, oficinas, pesquisas de campo, visitas técnicas, seminários e rodas de conversa. As intervenções podem assumir formatos como ações organizacionais, eventos, projetos de extensão, pesquisas aplicadas, planos de melhoria, desenvolvimento de produtos/serviços/processos, tecnologias sociais e campanhas educativas. São utilizados espaços como salas de aula, biblioteca, laboratório de informática, o Laboratório de Inovação, Empreendedorismo e Desenvolvimento Sustentável (LINEDS), além de organizações e espaços comunitários.
- **Avaliação:** Diagnóstica, processual e formativa, considerando a participação dos estudantes, o desenvolvimento das atividades propositivas, a capacidade de integração dos conhecimentos, a produção de sínteses individuais e coletivas, a comunicação dos resultados e a reflexão crítica sobre os processos. Instrumentos: relatórios, diários de bordo, portfólios, seminários, painéis, infográficos, mapas conceituais, estudos de caso, apresentações orais, autoavaliação e avaliação por pares.} \\ \hline
\multicolumn{3}{|p{\dimexpr\linewidth-2\tabcolsep-2\arrayrulewidth\relax}|}{\textcolor{ifscgreen}{\textbf{Bibliografia Básica:}}} \\ \hline
\multicolumn{3}{|p{\dimexpr\linewidth-2\tabcolsep-2\arrayrulewidth\relax}|}{1. FRIGOTTO, Gaudêncio; CIAVATTA, Maria; RAMOS, Marise Nogueira (org.). \textbf{Ensino médio integrado: concepção e contradições}. 3. ed. São Paulo: Cortez, 2012.
2. GIL, Antonio Carlos. \textbf{Como elaborar projetos de pesquisa}. 6. ed. São Paulo: Atlas, 2017.
3. MORIN, Edgar. \textbf{Os sete saberes necessários à educação do futuro}. Tradução de Catarina Eleonora F. da Silva, Jeanne Sawaya. 2. ed. rev. São Paulo: Cortez, 2011.
4. SILVA, Adriano Larentes da (org.) et al. \textbf{O currículo integrado no cotidiano da sala de aula}. Florianópolis: Publicação do IFSC, 2016. Disponível em: https://www.ifsc.edu.br/livros-e-periodicos.} \\ \hline
\multicolumn{3}{|p{\dimexpr\linewidth-2\tabcolsep-2\arrayrulewidth\relax}|}{\textcolor{ifscgreen}{\textbf{Bibliografia Complementar:}}} \\ \hline
\multicolumn{3}{|p{\dimexpr\linewidth-2\tabcolsep-2\arrayrulewidth\relax}|}{1. KEELLING, Ralph; BRANCO, Renato Henrique Ferreira. \textbf{Gestão de projetos: uma abordagem global}. Tradução de Cid Knipel Moreira. 4. ed. São Paulo: Saraiva, 2019.
2. MARCONI, Marina de Andrade; LAKATOS, Eva Maria. \textbf{Metodologia do trabalho científico: projetos de pesquisa, pesquisa bibliográfica, teses de doutorado, dissertações de mestrado, trabalhos de conclusão de curso}. 8. ed. São Paulo: Atlas, 2018.} \\ \hline
\end{xltabular}

\vspace{0.4cm}

